\documentclass[aspectratio=169,10pt]{beamer}

% ─── Thème et couleurs ───────────────────────────────────────────────────────
\usetheme{Madrid}
\definecolor{rtypeblue}{RGB}{0,80,150}
\definecolor{rtypedark}{RGB}{20,20,40}
\definecolor{rtypegreen}{RGB}{40,140,60}
\definecolor{rtypeorange}{RGB}{220,120,20}
\definecolor{rtypegray}{RGB}{100,100,100}
\definecolor{rtypelight}{RGB}{230,240,250}
\usecolortheme[named=rtypeblue]{structure}
\setbeamercolor{title}{fg=white,bg=rtypeblue}
\setbeamercolor{frametitle}{fg=white,bg=rtypeblue}
\setbeamercolor{block title}{bg=rtypeblue,fg=white}

% ─── Packages ────────────────────────────────────────────────────────────────
\usepackage[utf8]{inputenc}
\usepackage[T1]{fontenc}
\usepackage[french]{babel}
\usepackage{listings}
\usepackage{booktabs}
\usepackage{graphicx}
\usepackage{hyperref}
\usepackage{array}
\usepackage{adjustbox}
\usepackage{colortbl}
\arrayrulecolor{gray!50}
\renewcommand{\arraystretch}{1.3}
\usepackage{tikz}
\usetikzlibrary{positioning,arrows.meta,shapes.geometric,fit,calc,backgrounds}
\usepackage{forest}

% ─── Listings config ─────────────────────────────────────────────────────────
\lstset{
  language=C++,
  basicstyle=\ttfamily\tiny,
  keywordstyle=\color{rtypeblue}\bfseries,
  commentstyle=\color{gray},
  stringstyle=\color{red!60!black},
  breaklines=true,
  frame=single,
  columns=fullflexible,
  showstringspaces=false
}

% ─── Forest dirtree style ───────────────────────────────────────────────────
\forestset{
  dir tree/.style={
    for tree={
      font=\ttfamily\tiny,
      grow'=0,
      child anchor=west,
      parent anchor=south,
      anchor=west,
      calign=first,
      edge path={
        \noexpand\path [draw, -{Stealth[length=1.5pt]}]
        (!u.south west) +(4pt,0) |- node[fill,inner sep=0.6pt] {} (.child anchor)\forestoption{edge label};
      },
      before typesetting nodes={
        if n=1 {insert before={[,phantom]}} {}
      },
      fit=band,
      before computing xy={l=6pt},
      inner ysep=0pt,
      }
    }
}

% ─── Métadonnées ─────────────────────────────────────────────────────────────
\title[Bloc 3 — R-Type]{Bloc 3 — Créer une architecture logicielle pour applications lourdes}
\subtitle{Projet R-Type — Architecture et Intégration}
\author{Dytoma BATOGOUMA, Baptiste CAMERLYNCK, Jonathan UFARTE,\newline Noe ROBERTIES, Marc GRIOCHE}
\date{23 février 2026}

% ─── Navigation simplifiée ───────────────────────────────────────────────────
\setbeamertemplate{navigation symbols}{
  \insertframenavigationsymbol
}

% ─── Footline personnalisée ──────────────────────────────────────────────────
\setbeamertemplate{footline}{
  \leavevmode%
  \hbox{%
    \begin{beamercolorbox}[wd=.5\paperwidth,ht=2.25ex,dp=1ex,left]{author in head/foot}%
      \usebeamerfont{author in head/foot}\hspace*{1em}\insertshorttitle
    \end{beamercolorbox}%
    \begin{beamercolorbox}[wd=.4\paperwidth,ht=2.25ex,dp=1ex,center]{title in head/foot}%
      \usebeamerfont{title in head/foot}\insertdate
    \end{beamercolorbox}%
    \begin{beamercolorbox}[wd=.1\paperwidth,ht=2.25ex,dp=1ex,right]{date in head/foot}%
      \usebeamerfont{date in head/foot}\insertframenumber{} / \inserttotalframenumber\hspace*{1em}
    \end{beamercolorbox}%
  }%
  \vskip0pt%
}

% ══════════════════════════════════════════════════════════════════════════════
\begin{document}

% ─── PAGE DE TITRE ───────────────────────────────────────────────────────────
\begin{frame}
  \titlepage
\end{frame}

% ─── SOMMAIRE ────────────────────────────────────────────────────────────────
\begin{frame}{Sommaire}
  \tableofcontents
\end{frame}

% ══════════════════════════════════════════════════════════════════════════════
%  A7 — IMPLEMENTATION DES INTERACTIONS AVEC LE S.I.
% ══════════════════════════════════════════════════════════════════════════════
\section{A7 — Implémentation des interactions avec le S.I.}

% ──────────────────────────────────────────────────────────────────────────────
% SLIDE : B3-C15.1 — Interfaces (GUI/TUI)
% ──────────────────────────────────────────────────────────────────────────────
\begin{frame}[fragile]{B3-C15.1 — Interfaces (GUI)}
  \footnotesize
  L'application R-Type propose une interface graphique (GUI) riche gérée par \textbf{SDL2}, structurée en ``Pages'' pour optimiser le flux utilisateur.
  
  \begin{columns}[T]
  \begin{column}{0.48\textwidth}
  \textbf{Machine à états des menus (Menu.cpp) :}
  \begin{itemize}
    \item Gestion centralisée des états (\texttt{Page::HOME}, \texttt{Page::LOBBY}, etc.).
    \item Transition fluide sans rechargement complet de l'application.
    \item Composants réutilisables (\texttt{Button}, \texttt{InputBox}).
  \end{itemize}
  
  \vspace{0.2em}
  \begin{lstlisting}
  void Menu::activate(Registry& r, Page page, ...) {
      switch (page) {
      case Page::HOME: m_homePage.show(r); break;
      case Page::LOBBY: m_lobbyPage.show(r, ...); break;
      case Page::WIN: m_endPage.show(r, true); break;
      // ...
      }
      m_currentPage = page;
  }
  \end{lstlisting}
  \end{column}
  
  \begin{column}{0.48\textwidth}
  \textbf{Interactions utilisateur :}
  \begin{itemize}
    \item \textbf{Menus} : Navigation à la souris et au clavier.
    \item \textbf{Jeu} : Feedback visuel immédiat (animations, particules).
    \item \textbf{Lobby} : Mise à jour temps réel de la liste des joueurs et du chat.
  \end{itemize}
  
  \vspace{0.2em}
  \centering
  \begin{tikzpicture}[
    node distance=0.5cm,
    state/.style={draw, rounded corners, fill=rtypelight, font=\tiny\bfseries},
    arrow/.style={-{Stealth}, thick, rtypeblue}
  ]
    \node[state] (home) {HOME};
    \node[state, below left=of home] (conn) {CONNECTION};
    \node[state, below right=of home] (join) {JOIN LOBBY};
    \node[state, below=1.5cm of home] (lobby) {LOBBY};
    \node[state, below=of lobby] (game) {GAME};
    
    \draw[arrow] (home) -- (conn);
    \draw[arrow] (home) -- (join);
    \draw[arrow] (conn) |- (lobby);
    \draw[arrow] (join) |- (lobby);
    \draw[arrow] (lobby) -- (game);
    \draw[arrow] (game.east) to[out=0,in=0] (lobby.east);
  \end{tikzpicture}
  \end{column}
  \end{columns}
\end{frame}

% ──────────────────────────────────────────────────────────────────────────────
% SLIDE : B3-C15.2 — Accessibilité
% ──────────────────────────────────────────────────────────────────────────────
\begin{frame}{B3-C15.2 — Accessibilité et Ergonomie}
  \footnotesize
  L'interface est conçue pour être claire, réactive et navigable, respectant des critères d'accessibilité clés :
  
  \vspace{0.5em}
  \begin{columns}[T]
    \begin{column}{0.48\textwidth}
      \begin{block}{Ergonomie Visuelle}
        \begin{itemize}
          \item \textbf{Contraste élevé} : Textes clairs sur fonds sombres (thème spatial) pour une lisibilité maximale.
          \item \textbf{Police lisible} : Utilisation de polices sans-serif claires et de grande taille pour les éléments importants (scores, vies).
          \item \textbf{Feedback visuel} : Survol des boutons (hover), indicateurs de focus.
        \end{itemize}
      \end{block}
    \end{column}
    
    \begin{column}{0.48\textwidth}
      \begin{block}{Accessibilité des Contrôles}
        \begin{itemize}
          \item \textbf{Inputs Abstraits} : Le système \texttt{PlayerInputProcessor} mappe les actions (\texttt{GameInput::UP}, \texttt{ATTACK}) indépendamment du matériel.
          \item \textbf{Support Clavier/Souris} : Navigation complète possible dans les menus.
          \item \textbf{Messages d'erreur} : Retours explicites en cas d'échec de connexion ou de formulaire invalide.
        \end{itemize}
      \end{block}
    \end{column}
  \end{columns}
\end{frame}

% ──────────────────────────────────────────────────────────────────────────────
% SLIDE : B3-C16.1 — Intégrité des données (Implémentation)
% ──────────────────────────────────────────────────────────────────────────────
\begin{frame}[fragile]{B3-C16.1 — Intégrité des données (Implémentation)}
  \footnotesize
  L'intégrité des échanges réseau et des données critiques est assurée par notre protocole binaire custom sur UDP.
  
  \begin{columns}[T]
  \begin{column}{0.48\textwidth}
  \textbf{Structure des paquets (Message.hpp) :}
  \begin{itemize}
    \item \textbf{Sequence Number} : Détection des paquets perdus, dupliqués ou désordonnés.
    \item \textbf{Version} : Rejet des clients obsolètes.
    \item \textbf{Player ID} : Identification source (validée coté serveur).
  \end{itemize}
  
  \begin{lstlisting}
  class Message {
    uint16_t sequence_number; // Ordering
    uint32_t player_id;       // Source check
    uint8_t version;          // Compatibility
    std::vector<uint8_t> payload;
    
    // Serialisation binaire stricte
    std::vector<uint8_t> serialize() const;
  };
  \end{lstlisting}
  \end{column}
  
  \begin{column}{0.48\textwidth}
  \textbf{Validation à la réception (Server) :}
  \begin{itemize}
    \item Vérification de la taille du payload pour éviter les \textit{buffer overflows}.
    \item Validation que le \texttt{player\_id} correspond à l'IP source (prévention IP spoofing).
    \item Parsage strict via \texttt{Message::readBytes()} avec "bounds checking".
  \end{itemize}
  
  \vspace{0.3em}
  L'utilisation de \textbf{UDP} nécessite cette couche applicative pour garantir l'intégrité que TCP offrirait nativement, mais avec une latence bien moindre.
  \end{column}
  \end{columns}
\end{frame}

% ──────────────────────────────────────────────────────────────────────────────
% SLIDE : B3-C16.2 — Justification Intégrité/Confidentialité
% ──────────────────────────────────────────────────────────────────────────────
\begin{frame}{B3-C16.2 — Justification Intégrité et Confidentialité}
  \footnotesize
  
  \begin{block}{Confidentialité - Gestion des données sensibles (Server/DB)}
    \begin{itemize}
      \item \textbf{Mots de passe} : Hachage fort via \textbf{Argon2id} (bibliothèque \textit{libsodium}).
      \item \textbf{Stockage} : Base de données \textbf{SQLite3} locale au serveur, non exposée publiquement.
      \item \textbf{Principe} : Aucun mot de passe clair ne transite ou n'est stocké.
    \end{itemize}
  \end{block}
  
  \vspace{0.5em}
  \begin{block}{Intégrité - Stabilité de la solution}
    \begin{itemize}
      \item \textbf{Autorité Serveur} : Le client n'envoie que des intentions (inputs). Le serveur calcule et valide l'état du monde. Impossible pour un client de "téléporter" son vaisseau (intégrité logique du jeu).
      \item \textbf{Gestion mémoire} : Utilisation intensive de \texttt{std::unique\_ptr} et \texttt{RAII} pour empêcher les fuites de mémoire et la corruption de heap (\textit{memory corruption}).
      \item \textbf{Validation des Entrées} : Toutes les chaînes (usernames, chat) sont tronquées et désinfectées (taille max 255 chars).
    \end{itemize}
  \end{block}
\end{frame}

% ──────────────────────────────────────────────────────────────────────────────
% SLIDE : B3-C17.1 — Normes et Composants Tiers (Implémentation)
% ──────────────────────────────────────────────────────────────────────────────
\begin{frame}[fragile]{B3-C17.1 — Normes et Composants Tiers (Implémentation)}
  \footnotesize
  Le projet s'appuie sur des standards industriels et des bibliothèques robustes pour éviter de "réinventer la roue" sur les aspects critiques.
  
  \begin{columns}[T]
    \begin{column}{0.5\textwidth}
      \textbf{Standards utilisés :}
      \begin{itemize}
        \item \textbf{C++17} : Standard langage moderne (filesystem, optional).
        \item \textbf{CMake} : Standard de build system.
        \item \textbf{UDP/IPv4} : Protocole réseau standard (RFC 768).
        \item \textbf{JSON} : Format d'échange standard (RFC 8259).
      \end{itemize}
    \end{column}
    \begin{column}{0.5\textwidth}
      \textbf{Bibliothèques Tierces :}
      \begin{itemize}
        \item \textbf{SDL2} : Abstraction hardware (Audio/Vidéo/Input).
        \item \textbf{SQLite3} : Moteur SQL transactionnel ACID.
        \item \textbf{nlohmann/json} : Parsing JSON moderne en C++.
        \item \textbf{libsodium} : Cryptographie haute sécurité.
        \item \textbf{Prometheus-cpp} : Métriques et monitoring.
      \end{itemize}
    \end{column}
  \end{columns}
\end{frame}

% ──────────────────────────────────────────────────────────────────────────────
% SLIDE : B3-C17.2 — Justification des Choix Techniques
% ──────────────────────────────────────────────────────────────────────────────
\begin{frame}{B3-C17.2 — Justification des Choix Techniques}
  \footnotesize
  \begin{adjustbox}{max width=\textwidth}
  \begin{tabular}{l l l}
  \toprule
  \textbf{Composant} & \textbf{Rôle} & \textbf{Justification de la pertinence} \\
  \midrule
  \textbf{SDL2} & Rendu & \textit{Standard de facto} pour le développement 2D natif. Performance hardware, compatible Linux/Windows/Mac. \\ \hline
  \textbf{SQLite3} & Persistance & Zéro-configuration, simple fichier \texttt{.db}, suffisant pour nos besoins de matchmaking et auth. \\ \hline
  \textbf{UDP Binaire} & Réseau & Protocoles texte (HTTP/WS) trop lents pour le temps réel (60Hz). Binaire = 10x moins de bande passante. \\ \hline
  \textbf{CMake} & Build & Indispensable pour la portabilité (génération Makefile ou Visual Studio Solution). \\ \hline
  \textbf{Docker} & Déploiement & Garantit un environnement d'exécution identique (reproductibilité) pour le serveur. \\
  \bottomrule
  \end{tabular}
  \end{adjustbox}
  
  \vspace{0.5em}
  Ces choix garantissent \textbf{l'interopérabilité} (formats standards) et la \textbf{maintenance} (communautés actives pour chaque librairie).
\end{frame}

% ══════════════════════════════════════════════════════════════════════════════
%  A8 — MISE EN PLACE DE FONCTIONNALITÉS COMPLEXES
% ══════════════════════════════════════════════════════════════════════════════
\section{A8 — Fonctionnalités complexes et Code}

% ──────────────────────────────────────────────────────────────────────────────
% SLIDE : B3-C18.1 — Code Opérationnel et Fonctionnel
% ──────────────────────────────────────────────────────────────────────────────
\begin{frame}[fragile]{B3-C18.1 — Code Opérationnel (Complexité)}
  \scriptsize
  Le cœur du moteur de jeu (\texttt{GameInstance}) gère de manière autonome la simulation, intégrant physique, logique et réseau.
  
  \begin{columns}[T]
  \begin{column}{0.55\textwidth}
  \begin{lstlisting}[basicstyle=\ttfamily\tiny]
// GameInstance.cpp - Boucle de simulation
void GameInstance::update(double deltaTime) {
    // 1. Mise a jour des systemes ECS
    _registry.run_system<MovementSystem>(deltaTime);
    _registry.run_system<CollisionSystem>();
    _registry.run_system<WeaponSystem>();
    
    // 2. Gestion des inputs joueurs
    processInputs();
    
    // 3. Generation procedurale (Obstacles)
    _levelManager.update(deltaTime);
    
    // 4. Verification conditions de victoire
    checkWinCondition();
    
    // 5. Serialisation etat pour le reseau
    broadcastGameState();
}
  \end{lstlisting}
  \end{column}
  
  \begin{column}{0.4\textwidth}
  \textbf{Points forts du code :}
  \begin{itemize}
    \item \textbf{Autonomie} : Le serveur tourne en boucle infinie (\texttt{while(_running)}) et gère ses ticks (60 TPS).
    \item \textbf{Découplage} : La logique métier est séparée du réseau grâce à l'ECS.
    \item \textbf{Concurrence} : Utilisation de threads pour paralléliser I/O réseau et calculs physique (\texttt{LobbyManager}).
  \end{itemize}
  \end{column}
  \end{columns}
\end{frame}

% ──────────────────────────────────────────────────────────────────────────────
% SLIDE : B3-C18.2 — Conventions (Formatage et Nommage)
% ──────────────────────────────────────────────────────────────────────────────
\begin{frame}[fragile]{B3-C18.2 — Conventions de Formatage et Nommage}
  \footnotesize
  Le projet impose des conventions strictes pour garantir la lisibilité et la maintenance collaborative.
  
  \begin{columns}[T]
    \begin{column}{0.48\textwidth}
      \textbf{Outils Automatiques :}
      \begin{itemize}
        \item \textbf{.pre-commit-config.yaml} : Vérification automatique avant chaque commit.
        \item \textbf{Clang-Format} : Style \texttt{WebKit} imposé.
        \item \textbf{Standard} : C++17.
      \end{itemize}
      
      \vspace{0.2em}
      \begin{lstlisting}
# .pre-commit-config.yaml
repos:
  - repo: https://github.com/pre-commit/mirrors-clang-format
    hooks:
      - id: clang-format
        args: [--style=WebKit]
      \end{lstlisting}
    \end{column}
    
    \begin{column}{0.48\textwidth}
      \textbf{Conventions de Nommage :}
      \begin{itemize}
        \item \textbf{Classes} : \texttt{PascalCase} (ex: \texttt{GameInstance})
        \item \textbf{Méthodes} : \texttt{camelCase} (ex: \texttt{processInput})
        \item \textbf{Membres} : \texttt{m\_} ou \texttt{\_} prefix (ex: \texttt{\_registry}, \texttt{m\_currentPage})
        \item \textbf{Fichiers} : Identiques aux classes (ex: \texttt{Menu.cpp})
      \end{itemize}
      
      \vspace{0.3em}
      \textit{Respect strict vérifié lors des Code Reviews sur GitHub.}
    \end{column}
  \end{columns}
\end{frame}

% ──────────────────────────────────────────────────────────────────────────────
% SLIDE : B3-C19.1/19.2 — Intégration Tiers et Erreurs
% ──────────────────────────────────────────────────────────────────────────────
\begin{frame}[fragile]{B3-C19.1/19.2 — Intégration Tiers et Gestion d'Erreur}
  \footnotesize
  L'intégration des librairies externes se fait via des wrappers qui isolent la complexité et gèrent les erreurs de manière robuste.
  
  \begin{columns}[T]
  \begin{column}{0.48\textwidth}
  \textbf{Exemple : Initialisation DB (SQLite)}
  \begin{lstlisting}
void DB::init() {
    int rc = sqlite3_open(_path.c_str(), &_db);
    // Gestion d'erreur explicite (B3-C19.2)
    if (rc) {
        std::cerr << "Can't open database: " 
                  << sqlite3_errmsg(_db) << std::endl;
        // Fallback ou arret critique
        throw std::runtime_error("DB Error");
    }
    // ... creation tables
}
  \end{lstlisting}
  \end{column}
  
  \begin{column}{0.48\textwidth}
  \textbf{Exemple : Réseau (Sockets système)}
  \begin{lstlisting}
ssize_t UdpSocket::send(...) {
    ssize_t sent = sendto(_socket, ...);
    if (sent < 0) {
        // Enregistrement de l'erreur sans crasher
        perror("sendto failed");
        return -1;
    }
    return sent;
}
  \end{lstlisting}
  \end{column}
  \end{columns}
  
  \vspace{0.5em}
  La gestion des cas d'erreur tiers est encapsulée pour ne pas propager d'états invalides dans la logique métier (\texttt{GameInstance}).
\end{frame}

% ──────────────────────────────────────────────────────────────────────────────
% SLIDE DE FIN
% ──────────────────────────────────────────────────────────────────────────────
\begin{frame}
  \centering
  \vfill
  {\Huge \textbf{Merci de votre attention}}
  
  \vspace{1cm}
  \textbf{R-Type — Bloc 3}
  \vfill
\end{frame}

\end{document}
